% Figure: one gradient-descent step versus one Newton step on an
% anisotropic quadratic bowl f(x,y) = x^2/2 + 5 y^2/2, whose minimum is the
% origin. From the start (4,1) the negative gradient (-4,-5) points across
% the valley, so a gradient step overshoots the steep y-direction and
% undershoots the gentle x-direction; the Newton step -H^{-1} grad f lands
% on the minimum in one move because H^{-1} rescales each axis by its
% curvature.
\begin{figure}[htbp]
\centering
\begin{tikzpicture}[scale=1.0,
  contour/.style={enggray, thin},
  gdstep/.style={engred, very thick, -{Stealth[length=3mm]}},
  newton/.style={engblue, very thick, -{Stealth[length=3mm]}},
  pt/.style={circle, fill=engblack, inner sep=1.4pt},
]
  \begin{axis}[
    width=9.6cm, height=6.6cm,
    xlabel={$x$}, ylabel={$y$},
    xmin=-1.2, xmax=5, ymin=-1.6, ymax=1.8,
    axis lines=middle,
    xtick={-1,0,...,5}, ytick={-1,0,1},
    tick label style={font=\small},
    label style={font=\small},
    clip=false,
  ]
    % Elliptical contours of f = x^2/2 + 5 y^2/2 for a few levels.
    \foreach \c in {1,4,9} {
      \addplot[contour, domain=0:360, samples=80]
        ({sqrt(2*\c)*cos(x)}, {sqrt(2*\c/5)*sin(x)});
    }
    % Start point (4,1).
    \node[pt] at (axis cs:4,1) {};
    \node[anchor=south west, font=\small] at (axis cs:4,1) {start $(4,1)$};
    % Gradient descent step with eta = 0.18: grad = (4,5), step = (-0.72,-0.9)
    % lands at (3.28, 0.1). Draw the arrow.
    \draw[gdstep] (axis cs:4,1) -- (axis cs:3.28,0.1);
    \node[engred, anchor=west, font=\small] at (axis cs:3.35,0.55)
      {one GD step};
    % Newton step lands exactly at origin.
    \draw[newton] (axis cs:4,1) -- (axis cs:0,0);
    \node[engblue, anchor=south, font=\small] at (axis cs:2.0,0.62)
      {one Newton step};
    % Minimum marker.
    \node[pt] at (axis cs:0,0) {};
    \node[anchor=north east, font=\small] at (axis cs:0,-0.05) {minimum};
  \end{axis}
\end{tikzpicture}
\caption[One gradient step versus one Newton step on an anisotropic
bowl.]{The bowl $f(x, y) = \tfrac{1}{2}x^{2} + \tfrac{5}{2}y^{2}$ has
elliptical contours, five times steeper along $y$ than along $x$. From the
start $(4, 1)$ the negative gradient $(-4, -5)$ points mostly across the
valley, so a single gradient-descent step (red) drops the $y$-coordinate
almost to zero while barely reducing $x$, and the iterate must then crawl
along the $x$-axis. The Newton step $-H^{-1}\grad f$ (blue) premultiplies
the gradient by the inverse curvature and lands on the minimum in one move,
because $H^{-1}$ rescales each axis by its own steepness. The picture is the
geometry behind the condition-number remark: Newton's method works in the
coordinate system where every bowl is round.}
\label{fig:vol02:ch11:newton-step}
\end{figure}
